% Nouveau EGA V (cote 25): archive cote 25, site page 162 (PDF page 163).
% typescript. Draft compared: none.
% AI-transcribed by Claude Opus 5.5 (claude-opus-5-5, adaptive thinking, effort high), Message Batches API, batch ega25b (msgbatch_01A7n5YuXubUPNK1zsepi8RG), 2026-11-13; not yet checked by eye.
\scanpage{25/162}
\margin{\unsure{5}}

$\geqslant$ 2, on dira que x \hand{$\in$} X est un point singulier élémentaire de multiplicité
n, si son anneau local $\underline{O}_{X,x}$ est élémentaire de multiplicité n. (Signaler
en remarque que cette \xout{\illegible{1}} terminologie est en accord avec la notion
générale de multiplicité due à Samuel). Pour n=2, on parlera en particu-%
lier de singularité quadratique élémentaire (ou ordinaire, dans la termino
logie classique). On introduit de même, conformément à l'usage général,
les variantes ``géométriques'' pour X localement de type fini sur un corps:
x \inserted{$\in X$} est dit singularité géométriquement élémentaire de multiplicité n, si
pour toute (ou,ce qui revient au même\unsure{,} une)extension K de k, et tout
(ou un) point$^{z}$ de $X_{K}$ au dessus de x, rationnel sur K, z est une
singularité élémentaire de multiplicité n.

15.3. ``Généraliser''SUP prop.8 au cas de la multiplicité n.

15.4. ``Généraliser''SUP prop.9 et cor.10. au cas de la multiplicité n.
(Il ne s'agit pas vraiment de généralisations, mais de variantes..)

15.5. Introduire la notion de zéro singulier élémentaire de multiplicité
n (pour n=2, zéro singulier quadratique élémentaire, ou ordinaire) d'une
section $\emptyset$ de $\underline{O}_{X}$ , ou plus généralement d'un Module localament libre, sur
un X localement noethérien, et la notion géométrique correspondante
(sur un corps de base k).

15.6. Soit X un sous-préschéma d'un préschéma loc noeth régulier X', et
soit x \hand{$\in$}X. Pour que \unsure{\over{X}{x}|x} soit une singularité élémentaire de multiplicité
n de X, il f et s qu'il existe un\xout{e} voisinage ouvert U de x, et une
suite $\underline{O}_{U}$-régulière ($\emptyset_{1},..,\emptyset_{d}$) qui soit telle que $X|U = V(\emptyset_{1},..,\emptyset_{d})$ et
que x soit un zéro singulier élémentaire de multiplicité n de ($\emptyset_{1}$,..).
En particulier, si A est anneau local noethérien qui est élémentaire
\margin{$|$} de multiplicité n, alors A est un anneau ``intersection complète''; d'ail\xout{\illegible{1}}-%
leurs, on voit grâce à 15\unsure{.}4. que l'on peut trouver, au voisinage de x,

\mspage{162}
